\cotefile{1}
\scanpage{122}
\[
n!\alpha_n\leq\frac{\varepsilon_0}{n+1}\,\varepsilon_1\ldots\varepsilon_n\,e^n
\qquad \sqrt[n]{n!\alpha_n}\leq e\left(\frac{\varepsilon_0}{n+1}\right)^{1/n}
\deleted{\unsure{\varepsilon_1\ldots\varepsilon_n}}\unsure{(\varepsilon_1\ldots\varepsilon_n)^{1/n}|(\varepsilon_1\ldots\varepsilon_n)},
\]
\unsure{(} $n\to\infty$, $\left(\frac{\varepsilon_0}{n+1}\right)^{1/n}\to1$, et
$(\varepsilon_1\ldots\varepsilon_n)^{1/n}\to0$,
comme on voit immédiatement, d'où $\sqrt[n]{n!\alpha_n}\to0$,
d'où suit que $\unsure{S(\lambda)=}\sum n!\alpha_n(\lambda)$ \unsure{converge absolument} \unsure{à} \unsure{ch.} \unsure{en résulte|en vertu}
% slanted
\margin{par suite \unsure{convergence|converge} à \unsure{distance} \unsure{finie} sur $l^1$, la fonction $S(\lambda)$ est holomorphe}
Comme on a $|\alpha_n(\lambda)|\leq\|\lambda\|_1^n$, il suit que
la fonction $S(\lambda)$ \deleted{\unsure{sur $l^1$}} est holomorphe au voisinage
de l'origine\,; d'ailleurs, la forme \unsure{$n$ fois linéaire}
$\alpha_n(\lambda^{(1)},\ldots,\lambda^{(n)})$ associée à $\alpha_n(\lambda)$ est
\[
\alpha_n(\lambda^{(1)}\ldots\lambda^{(n)})=\frac{1}{n!}
\sum_{i_1,\ldots,i_n}\lambda^{(1)}_{i_1}\ldots\lambda^{(n)}_{i_n}
\]
(la somme étant étendue aux systèmes $(i_1,\ldots,i_n)$ d'indices \unsure{deux à deux|tous} distincts),
d'où $\|\alpha_n\|\leq\frac{1}{n!}$ \inserted{\deleted{\unsure{$\|\alpha_n\|$}} $n!\|\alpha_n\|\leq1$},
le rayon de convergence strict $R_0$ de la fonction $S(\lambda)$ à l'origine est
donc \deleted{\unsure{être|à peu}} $\geq1$. Soit alors $a_{\deleted{\unsure{m}}}^{(m)}$ l'élément de $l^1$
\deleted{\unsure{égal à}} donné par
\[
a_i^{(m)}=\begin{cases}\frac{1}{m}&\text{si }i\leq m\\0&\text{si }i>m\end{cases}
\]
on a manifestement $\|a^{(m)}\|_1=1$, et
\[
\alpha_n(a^m)=\deleted{\unsure{\frac{1}{m^n}}}
\left\{\begin{aligned}
&\deleted{\unsure{0\quad\text{si }n>m}}\\
&\frac{1}{m^n}\frac{m!}{n!(m-n)!}\quad\text{si }n\leq m
\end{aligned}\right.\,,\quad\text{d'où}
\]
\[
n!\alpha_n(a^m)=\frac{1}{m^n}\,m(m-1)\ldots(m-n+1)=
\deleted{\unsure{1}}(1-\tfrac{1}{m})(1-\tfrac{2}{m})\ldots(1-\tfrac{n-1}{m})
\]
Donc $n!\alpha_n(a^{(m)})\to1$ pour $m\to\infty$.
Il suit que $S(a^{(m)})\to\infty$ si $m\to\infty$,
donc $S(\lambda)$ n'est pas borné dans la boule
unité de $l^1$, et son rayon de convergence
large $R_1$ à l'origine est donc aussi $\leq1$, d'où $R_0=R_1=1$.

\deleted{Enfin, $S(\lambda)$ étant holomorphe au voisinage de \unsure{$0$|$a$} et pouvant être \unsure{prolongée|holomorphe|holomorphes} sur \unsure{chaque droite}}
% archive sheet number: 123
\scanpage{124}
\[
\bigl(\|u_1\wedge\cdots\wedge u_n\|_1\bigr)^2\leq
\unsure{\surd}\alpha_n(|u_1|\wedge\ldots\wedge|u_n|)
\,\alpha_n(|u_1^{\unsure{*|+}}|^{\prime}\!,\ldots\deleted{\unsure{\sqrt{\vphantom{u}|u_n|}}}|u_n|^{\prime})
\]
\[
|\alpha_n(u_1,\ldots u_n)|\leq\alpha_n(|u_1|,\ldots |u_n|).
\]

\uline{Démonstration} \quad Posons $u_i=\sum_\alpha\lambda_\alpha^{(i)}a_\alpha^{(i)}\otimes b_\alpha^{(i)}$
(représentation canonique de $u_i$) alors
\[
|u_i|=\sum_\alpha\lambda_\alpha^{(i)}a_\alpha^{(i)},\qquad
|u_i|'=\sum_\alpha\lambda_\alpha^{(i)}b_\alpha^{(i)}
\]
\[
u_1\wedge\cdots\wedge u_n=\frac{1}{n!}\sum_{\alpha_1\ldots\alpha_n}
\lambda_{\alpha_1}^{(1)}\ldots\lambda_{\alpha_n}^{(n)}
(a_{\alpha_1}^{(1)}\wedge\cdots\wedge a_{\alpha_n}^{(n)})
\otimes(b_{\alpha_1}^{(1)}\wedge\cdots\wedge b_{\alpha_n}^{\unsure{(n)}})
\]
d'où
\[
\|u_1\wedge\ldots u_n\|_1\leq\frac{1}{n!}\sum_{\alpha_1\ldots\alpha_n}
\lambda_{\alpha_1}^{(1)}\ldots\lambda_{\alpha_n}^{(n)}
\|a_{\alpha_1}^{(1)}\wedge\ldots\wedge a_{\alpha_n}^{(n)}\|
\|b_{\alpha_1}^{(1)}\wedge\ldots\wedge b_{\alpha_n}^{(n)}\|
\]
d'autre part on a, par \unsure{le $=$ calcul}
\[
\begin{aligned}
\alpha_n(|u_1|,\ldots |u_n|)&\inserted{=\operatorname{Tr}|u_1|\wedge\ldots |u_n|}=
\frac{1}{n!}\sum_{\alpha_1\ldots\alpha_n}\lambda_{\alpha_1}^{(1)}\ldots\lambda_{\alpha_n}^{(n)}
\|a_{\alpha_1}^{(1)}\wedge\ldots\wedge a_{\alpha_n}^{(n)}\|^2
\deleted{\unsure{b_{\alpha_1}^{(1)}\wedge\ldots\wedge a_{\alpha_n}^{(n)}}}\\
\alpha_n(|u_1|',\ldots |u_n|')&=\frac{1}{n!}\sum_{\alpha_1\ldots\alpha_n}
\lambda_{\alpha_1}^{(1)}\ldots\lambda_{\alpha_n}^{(n)}
\|b_{\alpha_1}^{1}\wedge\ldots\wedge b_{\alpha_n}^{(n)}\|^2
\end{aligned}
\]
\margin{$\alpha_n=$}
L'inégalité \unsure{voulue|cherchée} \deleted{\unsure{est}} \unsure{est|résulte} alors conséquence de l'inégalité
\deleted{$\unsure{(a_{\alpha_1}^{(1)}\wedge\ldots\wedge a_{\alpha_n}^{(n)},b_{\alpha_1}^{(1)}\wedge\ldots\wedge b_{\alpha_n}^{(n)})}$}
de Cauchy-Schwartz.

\fbox{Corollaire 1}
\[
|\alpha_n(u_1,\ldots u_n)|\leq\|u_1\wedge\ldots\wedge u_n\|_1
\leq\frac{1}{2}\bigl(\alpha_n(|u_1|,\ldots |u_n|)+\alpha_n(|u_1|',\ldots |u_n|')\bigr)
\]

\fbox{Corollaire 2}
\[
|\alpha_n(A+B)|^2\leq\deleted{\unsure{\sqrt{\vphantom A}}}\alpha_n(|A|+|B|)
\alpha_n(|A^*|'+|B^{\unsure{*}}|^{\unsure{+|'}})
\]
En effet, on a \deleted{\unsure{$\alpha_n(A+B$}}
\[
\begin{aligned}
|\alpha_n(A+B)|^{\unsure{2}}&=\left|\sum_k C_k^n\alpha_n(\overbrace{A,\ldots A}^{k},
\overbrace{B,\ldots B}^{n-k})\right|^{\unsure{2}}\\
&\leq\deleted{\unsure{\surd}}\biggl(\sum_k C_k^n
\bigl(\alpha_n(\overbrace{|A|,\ldots |A|}^{k},\overbrace{\unsure{|A||\lvert B\rvert},\ldots\unsure{B_1|\lvert B\rvert}}^{n-k})\bigr)^{1/2}
\bigl(\unsure{\sum_k}\alpha_n(\overbrace{|A|',\ldots |A|'}^{k},|B|'\ldots)\bigr)^{1/2}
\end{aligned}
\]
d'où, par l'inégalité de Cauchy-Schwartz
\[
|\alpha_n(A+B)|^{\unsure{2}}\leq\bigl(\sum_k C_k^n\alpha_n\unsure{(|A|,\ldots |A|,|B|,\ldots,|B|)}\bigr)^{1/2}
\bigl(\sum_k{}'\quad\bigr)^{\unsure{1/2}}\unsure{!}
\]
% archive sheet number: 125
\scanpage{125}
\[
=\bigl(\alpha_n(|A|+|B|)\bigr)^{\unsure{1/2}}\bigl(\alpha_n(|A|'+|B|')\bigr)^{1/2}
\]
\[
\deleted{\det(1+A+B)\leq\det(1+|A|)\det(1+|B|)}
\]

\fbox{Corollaire 3}
\[
|\det(1+A+B)|\leq\tfrac12\bigl(\det(1+|A|+|B|)+\det(1+|A^*|^{\unsure{\prime|,}}+|B^{\unsure{*}}|')\bigr)
\]
résulte aussitôt de
\[
|\alpha_n(A+B)|\leq\tfrac12\bigl[\alpha_n\unsure{(|A|+|B|)}+\alpha_n\unsure{(|A|'+|B|')}\bigr]
\]

\fbox{Corollaire 4}
\[
\boxed{|\det(1+A+B)|\leq\det(1+|A|)\det(1+|B|)}
\]
Il suffit de \deleted{\unsure{prouver}} le prouver pour $A$ et $B$
hermitiens positifs\,; car alors on aura, d'après
le corollaire 3,
\[
\begin{aligned}
|\det(1+A+B)|\leq\tfrac12\bigl[&\det(1+|A|)\det(1+|B|)\\
&+\det(1+|A|')\det(1+|B|')\bigr],
\end{aligned}
\]
ce qui donne l'inégalité voulue, en vertu de
$\det(1+|A|)=\det(1+\unsure{\lvert A(\,{}'|\lvert A\rvert'}),\ \det(1+|B|)=\det(1+|B|')$.

-- Quant au cas hermitien \inserted{positif}, on écrit
\[
\det(1+h+k)=\det(1+h)\bigl(1+\tfrac{\unsure{k|h|K}}{1+h}\bigr)=\det(1+h)\det(1+Hk)
\]
où $H=\frac{1}{1+h}$\,; on est ramené à prouver
$\det(1+Hk)\leq\det(1+k)$, il suffit pour cela
de prouver $\alpha_n(Hk)\leq\alpha_n(k)$ pour tout $n$,
soit $\operatorname{Tr}\overset{\unsure{n|\text{--}}}{\bigwedge}Hk\leq\|\overset{n}{\bigwedge}k\|_1$. Or on a
\[
\overset{n}{\bigwedge}Hk=(\overset{n}{\bigwedge}H)(\overset{n}{\bigwedge}k),\quad\text{d'où}
\]
\[
\|\overset{n}{\bigwedge}Hk\|_1\leq\|\overset{n}{\bigwedge}H\|_\infty\|\overset{n}{\bigwedge}k\|_1
\quad\text{et}\quad\|\overset{n}{\bigwedge}H\|_\infty\leq\|H\|_\infty^{\unsure{n|a}}
\]
\margin{\unsure{!}}or $\|H\|_\infty\leq1$.

\[
\boxed{\text{\fbox{Corollaire 5}}\quad\det(1+|A+B|)\leq\det(1+|A|)\det(1+|B|)}
\]
% archive sheet number: 126
\scanpage{126}
\[
=\bigl(\alpha_n(|A|+|B|)\bigr)^{\unsure{1/2}}\bigl(\alpha_n(|A|'+|B|')\bigr)^{1/2}
\]
\[
\deleted{\det(1+A+B)\leq\det(1+|A|)\det(1+|B|)}
\]

\fbox{Corollaire 3}
\[
|\det(1+A+B)|\leq\tfrac12\bigl(\det(1+|A|+|B|)+\det(1+|A^*|^{\unsure{\prime|,}}+|B^{\unsure{*}}|')\bigr)
\]
résulte aussitôt de
\[
|\alpha_n(A+B)|\leq\tfrac12\bigl[\alpha_n\unsure{(|A|+|B|)}+\alpha_n\unsure{(|A|'+|B|')}\bigr]
\]

\fbox{Corollaire 4}
\[
\boxed{|\det(1+A+B)|\leq\det(1+|A|)\det(1+|B|)}
\]
Il suffit de \deleted{\unsure{prouver}} le prouver pour $A$ et $B$
hermitiens positifs\,; car alors on aura, d'après
le corollaire 3,
\[
\begin{aligned}
|\det(1+A+B)|\leq\tfrac12\bigl[&\det(1+|A|)\det(1+|B|)\\
&+\det(1+|A|')\det(1+|B|')\bigr],
\end{aligned}
\]
ce qui donne l'inégalité voulue, en vertu de
$\det(1+|A|)=\det(1+\unsure{\lvert A(\,{}'|\lvert A\rvert'}),\ \det(1+|B|)=\det(1+|B|')$.

-- Quant au cas hermitien \inserted{positif}, on écrit
\[
\det(1+h+k)=\det(1+h)\bigl(1+\tfrac{\unsure{k|h|K}}{1+h}\bigr)=\det(1+h)\det(1+Hk)
\]
où $H=\frac{1}{1+h}$\,; on est ramené à prouver
$\det(1+Hk)\leq\det(1+k)$, il suffit pour cela

Ce corollaire contient le précédent, car on sait
que $\det(1+\deleted{\unsure{|}}A+B)\leq\det(1+|A+B|)$. -- Il
résulte du précédent, car on aura ``en général''
$|A+B|=U(A+B)$, $U$ unitaire, d'autre part
\[
\begin{aligned}
\det(1+UA+UB)&\leq\det(1+|UA|)\det(1+|UB|),\\
&=\det(1+|A|)\det(1+|B|)\qquad(\text{car } |UA|\\
&=U|A|U^{-1},\quad |UB|=U|B|U^{-1})
\end{aligned}
\]
% archive sheet number: 127
\scanpage{127}
\uline{Théorème} \unsure{Le déterminant} de Fredholm de $\unsure{u(\in L^1)|u\in L^1}$ \illegible{1}

Notation $\unsure{L_\infty|L_0|\mathcal L_0}$\,; dire \inserted{\unsure{est-ce que}} \unsure{exactement} quelles
parties du \deleted{\unsure{th}} théorème \unsure{--} sont connues.

\uline{Théorème}

1) Pour \deleted{\unsure{tout}} $0<p\leq\infty$, $\mathcal L^p(\mathcal H)$ est un idéal \inserted{bilatère}
dans $L(\mathcal H)$, stable pour l'involution $A\to A^*$. De
\deleted{\unsure{On a}} façon précise :
\[
\deleted{S_p}\|A\|_p=\deleted{S_p}\|A^*\|_p \tag*{(1)}
\]
\[
\begin{aligned}
\deleted{\unsure{S_p}}\|BA\|_p&\leq\|B\|\deleted{\unsure{S_p}}\|A\|_p
\qquad\text{pour } A\in\mathcal L^p(\mathcal H),B\in L(\mathcal H)\\
\|AB\|_p&\unsure{\leq|<}\|B\|\ \|A\|_p
\end{aligned}\tag*{(2)}
\]
Si $\infty\geq p\geq1$, $A\to\|A\|_p$ est une norme sur $\mathcal L^p(\mathcal H)$ \deleted{\unsure{(2)}}
i.e. on a
\[
\|A+B\|_p\leq\|A\|_p+\|B\|_p\quad\|\underset{\text{(1)}}{\lambda} A\|_p=|\lambda|\|A\|_p
\quad\unsure{(1\leq p\leq\infty}\tag*{(3)}
\]
\inserted{et $\mathcal L^p(\mathcal H)$ est complet pour cette norme}
Si $0<p\leq\infty$, alors l'expression $S_p(A)=\|A\|_p^p$ sur
$\mathcal L^p(A)$ satisfait à l'inégalité du triangle
\[
S_p(A+B)\leq S_p(A)+S_p(B)\qquad(0<p\leq1)\tag*{(4)}
\]
et muni de la distance, invariante par translation
$S_p(A-B)$, $\mathcal L^p(A)$ est un espace vectoriel topologique
métrisable et complet, non localement convexe
si $p<1$, et $\dim\mathcal H=+\infty$.
% slanted
\margin{\textcircled{2} \uline{Dual de $\mathcal L^p$}\,; \uline{réflexivité}}

\deleted{\unsure{2}}3) Si $A\in\mathcal L^p(\mathcal H)$ $(0<p<+\infty)$, alors la suite
des valeurs propres \inserted{$\alpha_i$} de $A$ est de puissance
$p$ième sommable, et on a
\[
\sum|\alpha_i|^p\leq S_p(A)\tag*{(5)}
\]
% slanted
\margin{généralisation aux normes de Schatten? \unsure{Ici plutôt|Ceci plutôt} en dernier \uline{Corollaire}}
$S_p(A)$ et $\|A\|_p$ sont des fonctions croissantes
sur l'ensemble des opérateurs hermitiens
positifs, i.e $H$ et $K$ hermitiens positifs, $H\leq K$,
implique $S_p(H)\leq S_p(K)$.

3) Soient $0<p,q,r\leq+\infty$, avec $\frac{1}{p}+\frac{1}{q}\geq\frac{1}{r}$. Si \deleted{donc}
$A\in\mathcal L^p(\mathcal H)$, $B\in\mathcal L^q(\mathcal H)$, alors $AB\in\mathcal L^r(\mathcal H)$, et on a
\[
\|AB\|_r\leq\|A\|_p\|B\|_q\tag*{(6)}
\]
\margin{(1) \ldots\ldots}
% archive sheet number: 128
\scanpage{128}
% slanted
\margin{Montrer qu'on peut se ramener à la dimension finie}
La première partie du théorème est résumée dans les
formules (1) à (4), \deleted{et} \inserted{outre} les affirmations relatives
aux \unsure{structures|topologies} vectorielles des espaces $\mathcal L^p(\mathcal H)$.
\inserted{Montrons \deleted{\unsure{d'abord}} d'abord que ces dernières sont \unsure{conséquences} très faciles \deleted{\unsure{de}} \unsure{des formules} \unsure{(3)|(1)} et (4)\ldots}
La formule (1) est évidente, car \deleted{\unsure{on sait que}}
\deleted{\unsure{bien} \unsure{au} cas $p<+\infty$, et} \deleted{\unsure{pour $p<+\infty$}} on sait (théorème
de Riesz) que $|A|$ et $|A^*|$ sont \deleted{\unsure{de même suite}} \inserted{unitairement équivalents}
\unsure{(?)|(2)} \deleted{valeurs propres}. La formule (2) est un cas
particulier de la formule (6), mais \deleted{\unsure{peut se}} \inserted{on peut}
donner une démonstration \deleted{\unsure{pourtant directe}}
\deleted{démontrer simplement de façon directe}
directe, en procédant comme dans $[\,]$.

En effet, si $(\alpha_i)$ \deleted{\unsure{resp. $(\beta_i)$}} désigne la
suite des valeurs \inserted{propres} de \deleted{$\lvert BA\rvert$} $|A|$
\inserted{et $(\beta_i)$ la suite des valeurs propres de $\unsure{\lvert BA\rvert|\lvert B\rvert A\rvert|\lvert BA\rvert^2}$ (\unsure{resp\textsuperscript{t}|resp.} $|AB|$)} \deleted{\unsure{resp. $|BA|$}},
rangées par ordre de grandeur
décroissante (chacune répétée autant de fois
que le comporte sa multiplicité) on a
\[
\beta_i\leq\|B\|\alpha_i\quad\text{pour tout }i.\quad(\text{loc cité, page \quad, lemme \quad}),
\]
d'où aussitôt les formules (2).

\deleted{\unsure{La dernière} des formules (3) est évidente, la}
Les formules (3) \deleted{peut} \inserted{peuvent} être regardées \deleted{de la formule} comme
\deleted{première est une conséquence immédiate}
conséquences de la formule (6), (que nous démontrerons de façon indépendante) de la façon
suivante. Si $p'$ est tel que $\frac{1}{p}+\frac{1}{p'}=1$,
et si $B\in\mathcal L^{p'}(\mathcal H)$, \deleted{et tel que}
\inserted{$\|B\|_{p'}\leq1$}, alors
on a en vertu de (6) $\|AB\|_1\leq\|A\|_p$, d'où
\[
|\operatorname{Tr}\unsure{(AB)|\lvert AB\rvert}|\leq\|A\|_p,\quad\text{et}\quad
\sup_{\|B\|_{p'}\leq1}|\operatorname{Tr}AB|\leq\|A\|_p
\]
% archive sheet number: 129
\scanpage{129}
Cette inégalité peut être remplacée par une égalité,
car il existe un \unsure{op\textsuperscript{r}|op.} \unsure{part\textsuperscript{t}} \unsure{is.|iso.} $U$ tel que
$UA=|A|$, \deleted{\unsure{puis, une}} et d'autre part, \unsure{comme|même} il
résulte aussitôt de la décomposition \unsure{spectrale|spectral} de
$|A|$ et du th. de Hölder \unsure{classique|classiques} sur les
espaces $L^p$ et $L^{p'}$, un opérateur hermitien
positif $H\in L^{p'}$ tel que $\operatorname{Tr}(H|A|)=\|A\|_p$,
soit \deleted{\unsure{$|$}}$\operatorname{Tr}(HUA)$\deleted{\unsure{$|$}}$=\|A\|_p$. Si on pose
$B=UA$, on aura $\|B\|_{p'}\leq\|U\|\ \|H\|_{p'}$ (formule
(1)) donc $\|B\|_{p'}\leq1$, ce qui prouve notre assertion.
Par suite on a bien
\[
\|A\|_p=\sup_{\|B\|_{p'}\leq1}|\operatorname{Tr}AB|\tag*{(7)}
\]
ce qui prouve que $\|A\|_p$ est bien une
\unsure{semi norme|seminorme|semi-norme}, comme borne supérieure de
semi-normes $A\to|\operatorname{Tr}AB|$

Une telle démonstration n'est plus \unsure{pratiquable|praticable}
pour $p<1$\,; \deleted{et} nous démontrerons la formule (4)
par voie analytique, grace au

\uline{Lemme} 1 \quad Soit $(\alpha_i)$ une suite \deleted{\unsure{de puiss}} \deleted{décroissante} de
nombres $\geq0$, \deleted{\unsure{de puissance $p$ième sommables}}
\inserted{et soit} \deleted{\unsure{Soit|soit}} $0<p<1$. Posons $f(z)=\prod_{i=1}^{\infty}(1+\alpha_i z)$.
Alors on a
\[
\sum\alpha_i^p=p\frac{\sin\pi p}{\pi}\int_0^\infty\frac{\log f(r)}{r^{1+p}}\,dr\tag*{(8)}
\]
(les deux membres pouvant éventuellement
être infinis)
% archive sheet number: 130
\scanpage{130}
\uline{Corollaire}\quad Si $A\in\mathcal L^1(\mathcal H)$, soit
$M_A(r)=\sup_{|z|=r}\det(1+zA)$
\deleted{si $A$ est hermitien positif, et si $A\in\mathcal L^p(\mathcal H)$}, \deleted{\unsure{tel|on}}
\deleted{\unsure{avec}} \inserted{et} $0<p<1$\,; on a
\[
S_p(A)=p\frac{\sin\pi p}{\pi}\int_0^\infty\frac{\log M_{|A|}(r)}{r^{1+p}}\,dr\tag*{(9)}
\]

Il suffit en effet de poser $f(z)=\det(1+z|A|)$,
les $\alpha_i$ étant les valeurs propres de $|A|$\,; on aura
ici $M(r)=f(r)$.

Pour démontrer (4), on \deleted{\unsure{pe|peut}} peut écrire
\[
S_p(A+B)=S_p(|A+B|)=p\frac{\sin\pi p}{\pi}\int_0^\infty
\frac{\log M_{|A+B|}(r)}{r^{1+p}}\,dr
\]
Or $M_{|A+B|}(r)=\det(1+\unsure{r|v}|A+B|)\leq\det(1+r|A|)\det(1+\unsure{r|v}|B|)$
(formule \unsure{\ldots}) \deleted{\unsure{est donc|soit $M_A$}} soit
\[
M_{|A+B|}(r)\leq M_{|A|}(r)\unsure{\cdot|+}M_{|B|}(r)\tag*{(10)}
\]
d'où
\[
S_p(A+B)\leq\frac{p\sin\pi p}{\pi}\int_0^\infty\frac{\log M_{|A|}(r)}{r^{1+\unsure{p|r}}}\,dr
+\frac{p\sin\pi p}{\pi}\int_0^\infty\frac{\log M_{|B|}(r)}{r^{1+\unsure{p|r}}}\,dr
\]
Appliquant de nouveau (9), on obtient bien la
formule voulue (4).

\deleted{Pour prouver}

2) Pour prouver la formule (5), je ne connais
qu'une méthode analytique. De façon générale,
\unsure{.} si $A$ est un opérateur cp\textsuperscript{t} continu, \unsure{soient|on prend|soit} la
suite des val. propres $(\alpha_i)$, \deleted{\unsure{prises|posons}} et si
$0<p<+\infty$, posons $s_p(A)=\sum|\alpha_i|^p$.
\deleted{On a $\unsure{s_p(A)|S_p(A)}=$ \unsure{Si|$S_p$} évidemment pour tout entier $n>0$ $s_p(A)=s_{p/n}(A^n)$. Supposons}
% archive sheet number: 131
\scanpage{131}
Supposons d'abord $p<1$, alors on a, d'après
une formule de majoration classique pour la
somme des puissances $p$\textsuperscript{ièmes} des inverses des
modules des zéros d'une fonction entière
\[
f(z)=\det(1\unsure{-|+}zA)\qquad(0<p<1)
\]
\[
s_p(A)\leq p^2\int_0^\infty\frac{\log M_A(r)}{r^{p+1}}\,dr\tag*{\deleted{(10)}}
\]
or $M_A(r)\leq M_{|A|}(r)$ (\unsure{\ldots\ --}) d'où
\[
s_p(A)\leq p^2\int_0^\infty\frac{\log M_{|A|}(r)}{r^{\unsure{p-1|p+1}}}\,dr
\]
\,; appliquant alors la formule
(9), on obtient donc la majoration
\[
s_p(A)\leq\frac{p\pi}{\sin p\pi}S_p(A)\tag*{(10)}
\]
Le facteur $\frac{p\pi}{\sin p\pi}$ est $>1$, mais \unsure{tend} vers 1
pour $p\to0$. Soit alors $0<p\unsure{<|\leq}+\infty$ quelconque,
et soit $n$ un entier $>0$ tel que $\frac{p}{n}<1$.
On a évidemment $s_p(A)=s_{p/n}(A^n)$, d'où, en
vertu de (10) :
\[
s_p(A)\leq\frac{\frac{p}{n}\pi}{\sin\frac{p}{n}\pi}S_{p/n}(A^n)
\]
Or, on a \deleted{\unsure{$S_{p/n}(A^n)\leq S_p(A)$}}
\[
\|A^n\|_{p/n}\leq\underbrace{\|A\|_p\ldots\|A\|_p}_{n\ \text{facteurs}},
\]
comme il résulte aussitôt de la formule (6), qui sera
démontrée plus bas (alors il est immédiat
que cette formule se généralise au produit d'un
nb fini quelconque d'opérateurs), d'où
\[
S_{p/n}(A^{\unsure{n|{-}}})=\left(\|A^n\|_{p/n}\right)^{p/n}
\leq\|A\|_p^p=S_p(A).
\]
% archive sheet number: 132
\scanpage{132}
d'où
\[
s_p(A)\leq\frac{\pi\frac{p}{n}}{\sin\pi\frac{p}{n}}S_p(A).
\]
Faisant tendre $n$ vers l'infini, on trouve l'inégalité (5).

L'assertion que si $0\leq H\leq K$, \deleted{sont hermitiens positifs} \inserted{alors on a}
$S_p(H)\leq S_p(K)$ résulte aussitôt du fait
bien connu que sous les conditions envisagées,
la suite des valeurs propres \inserted{décroissante} de $H$ est majorée
par la suite des valeurs propres décroissantes
de $K$ \deleted{\unsure{C}} (comme il résulte de la caractérisation extrémale \inserted{classique} \deleted{\unsure{des}} de la suite des valeurs
propres d'un opérateur hermitien complètement continu).

\deleted{3) Pour prouver la formule (6), on \unsure{écr} \inserted{écrit}
\unsure{écrira|écrire}, pour un entier $m>0$ arbitraire}
\[
\deleted{S_r(AB)=S_{r/2}\bigl(AB(AB)^*\bigr)=
\unsure{S_{r/2m}\bigl((AB(AB)^*)^m\bigr)}}\quad
\deleted{\inserted{S_{r/2}(\unsure{L|H})\quad\text{où }L=(AB)(AB)^*.}}
\]
\deleted{Si $m$ est un entier $>0$, on aura aussi}
\[
\deleted{S_r(AB)=S_{r/2}\,L=S_{r/2m}(L^m),\quad\text{d'où, en vertu de la formule (9),}}
\]
\deleted{en posant \unsure{$k(r)=\frac{\pi r}{\sin\pi r}$} \unsure{$k(r')$} \unsure{$\frac{1}{k(r')}$}}
\margin{\deleted{$k(\unsure{u|x|\alpha})=\frac{\pi u}{\sin\pi u}$}}
\deleted{et $S_r(AB)$ \unsure{et si $k$} $r'=\frac{r}{2m}$}
\[
\deleted{S_r(AB)=\frac{\unsure{r'}}{k(\unsure{r/2|r'})}\int_0^\infty
\frac{\log M_{L^m}(\rho)}{\rho^{1+r/2m}}\,\unsure{dr|d\rho|dv}}
\]
\deleted{Posons $r'=\frac{r}{2m}$\quad $\unsure{r'|p'}=\frac{p}{2m}$\quad $q'=\frac{q}{2m}$, on aura
évidemment encore $\frac{1}{p'}+\frac{1}{q'}\geq\frac{1}{r'}$. La formule (9)
donne alors}
\[
\deleted{S_r(AB)=S_{r'}(L^m)=\frac{r'\sin\pi r'}{\unsure{\pi|k(r')}}
\int_0^\infty\frac{\log M_{L^m}(\rho)}{\rho^{1+r'}}\,\unsure{dr|d\rho|dv}.}
\]
% archive sheet number: 133
\scanpage{133}
3) La preuve de la formule (6) est la plus difficile.
Nous commençons par donner une
preuve directe et très simple pour le cas où
\deleted{$\unsure{p,q,r\geq1,\ r=1}$} \inserted{$p\geq1,q\geq1,r=1$} *
(preuve qui a d'ailleurs l'avantage
de \inserted{devoir} s'appliquer sans modification au cas d'une
trace \unsure{$\tau$|$\sigma$|$\operatorname{Tr}$|$D$} au sens de von Neumann - Dixmier, définie
sur une algèbre auto-adjointe $\unsure{\mathcal A|\mathfrak A}$ \deleted{\unsure{ft fermée}}
\deleted{\unsure{unift fermée}} \unsure{unif\textsuperscript{t}|unift} fermée \unsure{fortement}
(1) (\deleted{\unsure{les éléments de}} \unsure{(1)} les ``\deleted{éléments} \inserted{opérateurs} de puissance
$p$\textsuperscript{ième} sommable'' de $\mathcal A$ \inserted{au sens de la trace $\tau$}
se définissant encore
comme l'ens des $A\in\mathcal A$ tels que
$\tau(|A|^p)<+\infty$.)\unsure{\textsuperscript{*}}.
% slanted
\margin{?? Il n'est d'ailleurs pas difficile \deleted{\unsure{de}} \unsure{étendre|d'étendre} alors le même \deleted{\unsure{tout}} énoncé au cas où $r\geq1$}
(Au N° suivant, nous
obtiendrons d'ailleurs le $=$ résultat comme
conséquence d'une proposition plus générale).

Nous commençons par prouver que, si
$A_1,\ldots A_{2^n}$ sont $2^n$ opérateurs cpt continus,
($n$ entier $\geq1$) alors on a \textsuperscript{(1)}
\[
\|A_1 A_2\ldots A_{2^n}\|_1\leq\|A_1\|_{2^n}\,\|A_2\|_{2^n}\ldots\|A_{2^n}\|_{2^n}\tag*{(11)}
\]
Nous le prouvons par récurrence sur $n$\,;
c'est vrai si $n=1$, \deleted{\unsure{on a bien}} \inserted{(formule \unsure{\ldots})}, supposons
\deleted{\unsure{$\|A_1A_2\|_1\leq\|A_1\|_2$|$\|A_1A_2\|_1\leq\|A_1\|_2\|A_2\|_2$}}
la \inserted{formule analogue} démontrée pour $2^{n-1}$ facteurs, et démontrons
alors (11). Écrivant $A_1\ldots A_{2^n}=(A_1A_2)\ldots(A_{2^n-1}A_{2^n})$
le premier membre sera, en vertu de l'hyp.
de récurrence, majoré par $\|A_1A_2\|_{2^{n-1}}\ldots\|A_{2^n-1}A_{2^n}\|_{2^{\unsure{n-1}}}$

\noindent\rule{\linewidth}{0.3pt}

(1) Le lecteur notera que ce résultat suffirait pour
démontrer \deleted{directement} \inserted{directement}, comme nous le ferons plus bas,
la formule (6) dans le cas général.
% archive sheet number: 134
\scanpage{134}
de sorte qu'on est ramené à démontrer des formules telles que
$\|A_1 A_2\|_{2^{n-1}}\leq\|A_1\deleted{\unsure{\|_{2^2}}}\|_{2^n}\|\unsure{A_2|B_1}\|_{2^n}$, soit, \unsure{précisément|généralement},
\uline{$\|AB\|_{2^{n-1}}\unsure{<|\leq}\|A\|_{2^{\unsure{n|n-1}}}\,\|B\|_{2^n}$}.
Or le premier membre s'écrit
\[
\|AB\|_{2^{n-1}}=\left(\bigl\|\bigl(AB(AB)^*\bigr)^{\frac{2^{n-1}}{2}}\bigr\|_1\right)^{\frac1{2^{n-1}}}
=\left(\|C\|_1\right)^{\frac1{2^{n-1}}}
\]
où
\[
\unsure{C|C_1}=(ABB^*A^*)^{2^{n-2}}=A\,(BB^*)\,(A^*A)\ldots(BB^*)\,A^*
\]
($2^{n-1}+1$ facteurs au dernier membre) on a
\[
\|C\|_1=\operatorname{Tr} C\deleted{\unsure{=\operatorname{Tr}}}\inserted{=\operatorname{Tr}\deleted{\unsure{C}}C'}\quad
\inserted{\text{où }C'=}\, A^*A\,.\,BB^*\,.\,A^*A\,.\,\ldots\,BB^*\deleted{\unsure{)=\operatorname{Tr}\lvert C\rvert}}
\]
\deleted{où dans le dernier membre figurent $2^{n-1}$ facteurs}
($2^{n-1}$ facteurs au second membre)

Or $|\operatorname{Tr} C'|\leq\|C'\|_1$, et d'après l'hypothèse de récurrence, on a de nouveau
\[
\begin{aligned}
\|C'\|_1&\leq\|A^*A\|_{2^{n-1}}\|BB^*\|_{2^{n-1}}\ldots\|BB^*\|_{2^{n-1}}\\
&=\bigl(\|A^*A\|_{2^{n-1}}\bigr)^{2^{n-2}}\bigl(\|B^*B\|_{2^{n-1}}\bigr)^{2^{n-2}}
\end{aligned}
\]
Or $\|A^*A\|_{2^{n-1}}=\bigl(\|A\deleted{\unsure{_2}}\|_{2^n}\bigr)^2$,
$\|B^*B\|_{2^{n-1}}=\bigl(\|B\|_{2^n}\bigr)^2$
d'où $\|C'\|_1\leq\bigl(\|A\|_{2^n}\|B\|_{2^n}\bigr)^{2^{n-1}}$. On obtient donc bien
\[
\|AB\|_{2^{n-1}}=\bigl(\|C\|_1\bigr)^{\frac1{2^{n-1}}}
\leq\bigl(\|C'\|_1\bigr)^{\frac1{2^{n-1}}}\leq\|A\|_{2^n}\unsure{\lvert B\||\|B\|}_{2^n},
\]
ce qui prouve la formule (11)

\medskip\noindent\hspace*{2em}-\enspace-\enspace-
% archive sheet number: 135
\scanpage{135}
\deleted{\unsure{2)|4)}} Soient $p\geq\unsure{1|0}$,
$\unsure{q\geq0|q\leq0|q\geq1}$ \quad $\frac1p+\frac1q=1$,
$A\in\mathcal L^p$, $B\in\mathcal L^q$, alors $AB\in\mathcal L^1$, et
$\|AB\|_1\leq\|A\|_p\,\|B\|_q$

\uline{Preuve} On peut supposer $A,B\in\mathcal L^1$
On a $A=UH$, $B=KV$ ($U,V$ partiellement isométriques)

Pour \deleted{\unsure{tout}} $n$, soient $\alpha,\beta$ deux entiers tels que $\alpha+\beta=2^n$;
posons $h=H^{\frac1\alpha}$, $k=K^{\frac1\beta}$, \ldots
\[
AB=Uh^\alpha k^\beta V\qquad\text{d'où}\quad\|AB\|_1\leq\|h^\alpha k^\beta\|_1\quad\unsure{<|\leq}
\]
or (lemme 1)
\[
\begin{aligned}
\|h^\alpha k^\beta\|_1&\leq
\underbrace{\|h\|_{2^n}\ldots\|h\|_{2^n}}_{\alpha}
\underbrace{\|k\|_{2^n}\ldots\|\unsure{k|h}\|_{2^n}}_{\beta}\\
&=\bigl(\operatorname{Tr} h^{2^n}\bigr)^{\frac\alpha{2^n}}
\bigl(\operatorname{Tr}(k^{2^n})\bigr)^{\frac\beta{2^n}}\qquad\text{d'où}
\end{aligned}
\]
\[
\|AB\|_1\leq\bigl(\operatorname{Tr}H^{\deleted{\unsure{2|\frac{2^n}{\alpha}}}\lambda}\bigr)^{\frac1\lambda}
\bigl(\operatorname{Tr}K^\mu\bigr)^{\frac1{\unsure{r|\mu}}}
\quad\text{où }\lambda=\frac{2^n}{\alpha},\ \mu=\frac{2^n}{\beta}
\]
\deleted{Faisons alors}

Cette inégalité est vraie pour \inserted{\unsure{un couple|on voit}} $(\lambda,\mu>1$ pourvu qu'on puisse les mettre sous la \deleted{\unsure{forme}} forme
$\frac{2^n}{\alpha},\frac{2^n}{\beta}$ avec $\alpha+\beta=2^n$,
i.e pourvu $\frac1\lambda+\frac1\mu=1$, \deleted{\unsure{que}}
et $\frac1\lambda$ et $\frac1\mu$ sont des multiples de $\frac1{2^n}$;
par un passage à la limite \unsure{immédiat|immédiate|évident}, on conclut
\[
\|AB\|_1\leq\bigl(\operatorname{Tr}H^p\bigr)^{\frac1p}
\bigl(\operatorname{Tr}H^{\unsure{p'|q'}}\bigr)^{\frac1{p'}}
\]
\inserted{On démontre de même, si $\frac1{p_1}+\frac1{p_2}\ldots+\frac1{p_n}=1$,
$\|A_1\ldots A_n\|_1\leq\|A_1\|_{p_1}\ldots\|A_n\|_{p_n}$}

3) On en conclut qu'on a ($p>1,\ \frac1p+\frac1{p'}\unsure{=|\geq}1$)
\[
\|A\|_p=\sup_{\|B\|_{\unsure{q|p'}}\leq1}|\operatorname{Tr} AB|\qquad\text{(d'où facilement}
\]
\[
\|A_1+A_2\|_p\leq\|A_1\|_p+\|A_2\|_p\quad)\quad\text{et enfin}
\]
\deleted{$AB$} si $\frac1p+\frac1q\geq\frac1r$ \quad ($p\geq1,q\geq1,\unsure{r}\geq1$)
% archive sheet number: 136
\scanpage{136}
Le résultat précédent nous permet de démontrer le résultat auxiliaire suivant, qui nous permettra de démontrer la formule (6) en \deleted{toute} généralité :

\uline{Lemme}. Soient $A$ et $B$ deux opérateurs $\in\mathcal L^{2m}$ \unsure{($\mathcal H$)|($\mathfrak H$)|($H$)}, posons
$H=A^*A,\ K=B^*B,\ L=(AB)^*(AB)$.
\deleted{\unsure{$L$ soit|et soit}} \deleted{\unsure{$L^m$ est, $H$|$L^n,H^n$}} Soient $m$ et $n$ des entiers $\geq1$.
Alors \deleted{\unsure{$L^m$,|$L^n$,}} $L^m,H^m$ et $K^m$ sont opérateurs à trace, et on a
\[
\alpha_n\bigl(|\deleted{\unsure{A}}L|^{\unsure{m}}\bigr)
\leq\alpha_n\bigl(\deleted{\unsure{H}}H^m\bigr)\,\alpha_n\bigl(\deleted{\unsure{K}}K^m\bigr)
\tag*{(12)}
\]
Le fait que $L^m,H^m,K^m$ soient opérateurs à trace résulte \deleted{enfin d'un des précédents} \inserted{aussitôt de ce qui précède}. (12) s'écrit alors, \deleted{\unsure{cela|et}}
\deleted{\unsure{par contre s'écrit en effet en posant}}
\inserted{en posant} $\mathscr A=\underline{\wedge} A,\ \mathscr B=\underline{\wedge} B$
\inserted{$\mathscr H=\mathscr A^*\mathscr A,\ \mathscr K=\mathscr B^*\mathscr B,\ \mathscr L=(\mathscr A\mathscr B)^*(\mathscr A\mathscr B)$}
\[
\operatorname{Tr}\,|\deleted{\unsure{\mathscr A\mathscr B}}\mathscr L|^{\unsure{m}}\leq
\operatorname{Tr}\,\deleted{\unsure{\mathscr H}}\mathscr H^m\ \operatorname{Tr}\,\deleted{\unsure{\mathscr B}}\mathscr K^m,
\]
\deleted{\illegible{2}}
\[
\deleted{\unsure{\bigl(\|\mathscr A\mathscr B\|_m\bigr)^m\leq
\bigl(\|\mathscr A\|_m\bigr)^m\bigl(\|\mathscr B\|_m\bigr)^m|\bigl(\|L\|_m\bigr)^m\leq\bigl(\|H\|_m\bigr)^m\bigl(\|K\|_m\bigr)^m}}\quad
\deleted{\text{\unsure{ce qui}}}
\]
\deleted{$\|\mathscr A\mathscr B\|_m\leq\|\mathscr A\|_m\|\mathscr B\|_m$. Or, on a vu qu'on a même
$\|AB\|_m\unsure{<|\leq}\|A\|_{\frac m2}\|B\|_{\frac m2}$}

ce qui nous ramène au cas où $n=1$. On doit alors démontrer
\[
\unsure{\operatorname{Tr}\||\operatorname{Tr}\lvert|\operatorname{Tr}[}\,(ABB^*A^*)\ldots(ABB^*A^*)\unsure{\|_1|\rvert|]_1|]_m}
\leq\bigl(\|\deleted{\unsure{H}}A^*A\|_m\bigr)^m\bigl(\|B^*B\|_m\bigr)^m
\]
Or, le premier membre \deleted{est la trace-norme} est aussi égal à la
\deleted{d'un produit de $4m$ facteurs, est en effet} trace du produit de $2m$ facteurs
\deleted{majoré par}
\[
(A^*A)(BB^*)\ldots(BB^*),
\]
donc au plus égal à la trace-norme de ce produit, donc au plus égal à
$\bigl(\|A^*A\|_{2m}\bigr)^m\bigl(\|BB^*\|_{2m}\bigr)^m$
d'après la formule $\|A_1\ldots A_{2m}\|_1\leq\|A_1\|_{2m}\ldots\|A_{2m}\|_{\unsure{2m}}$

Or, on a $\|A^*A\|_{2m}\leq\|A^*A\|_m$, et
$\|BB^*\|_{2m}\unsure{=|\leq}\|B^*B\|_{2m}\leq\|B^*B\|_m$, d'où résulte le lemme.
% archive sheet number: 137
